\documentclass[11pt,a4paper]{article}
\usepackage[utf8]{inputenc}
\usepackage[T1]{fontenc}
\usepackage{lmodern}
\usepackage[french]{babel}
\usepackage{amsmath,amssymb,mathtools}
\usepackage{icomma}
\usepackage[version=4]{mhchem}
\usepackage{siunitx}
\sisetup{output-decimal-marker={,},per-mode=symbol}
\usepackage{xcolor}
\usepackage{colortbl}
\usepackage{tikz}
\usepackage[most]{tcolorbox}
\usepackage{enumitem}
\usepackage{array}
\usepackage{booktabs}
\usepackage{fancyhdr}
\usepackage[hidelinks]{hyperref}
\usetikzlibrary{arrows.meta,positioning,calc}
\usepackage[a4paper,margin=2.1cm,headheight=26pt,headsep=10pt]{geometry}

\definecolor{primary}{RGB}{146,95,161}
\definecolor{primarydark}{RGB}{92,55,108}
\definecolor{acidecol}{RGB}{205,60,45}
\definecolor{basecol}{RGB}{40,110,200}
\definecolor{excol}{RGB}{0,135,90}
\definecolor{attncol}{RGB}{200,40,55}

\tcbset{boxbase/.style={enhanced,breakable,boxrule=0.9pt,arc=2.5pt,left=3mm,right=3mm,top=2.3mm,bottom=2.3mm,fonttitle=\bfseries,coltitle=white,toptitle=0.6mm,bottomtitle=0.6mm}}
\newtcolorbox[auto counter]{corrbox}[1][]{boxbase,colback=excol!4,colframe=excol,title={\textsc{Corrigé}~\thetcbcounter\ifx\relax#1\relax\relax\else\textmd{ --- \itshape #1}\fi}}
\newtcolorbox{astuce}{boxbase,colback=attncol!6,colframe=attncol,sharp corners}

\pagestyle{fancy}
\fancyhf{}
\fancyhead[L]{\small\textcolor{primarydark}{\textbf{Fiche 10} — Thème 1 : couples acide/base}}
\fancyhead[R]{\small\textcolor{primarydark}{\textsc{Corrigé — Difficile}}}
\fancyfoot[C]{\small\thepage}
\renewcommand{\headrulewidth}{0.8pt}
\renewcommand{\headrule}{\hbox to\headwidth{\color{primary}\leaders\hrule height \headrulewidth\hfill}}

\setlength{\parindent}{0pt}
\setlength{\parskip}{3pt}

\begin{document}

\begin{center}
{\LARGE\bfseries\color{primarydark} Thème 1 --- Corrigé Difficile}
\end{center}

\begin{corrbox}[l'acide phosphorique, un triacide]
\begin{enumerate}[label=\alph*),leftmargin=1.6em,itemsep=3pt]
  \item Les trois pertes successives :
  \[ \ce{H3PO4 <=> H2PO4- + H+} \quad;\quad \ce{H2PO4- <=> HPO4^{2-} + H+} \quad;\quad \ce{HPO4^{2-} <=> PO4^{3-} + H+}. \]
  D'où les trois couples : \ce{H3PO4}/\ce{H2PO4-}, \ce{H2PO4-}/\ce{HPO4^{2-}}, \ce{HPO4^{2-}}/\ce{PO4^{3-}}.
  \item Ampholytes : \textbf{\ce{H2PO4-}} et \textbf{\ce{HPO4^{2-}}}. Chacune apparaît une fois comme acide et
        une fois comme base dans la liste des couples : elle possède encore un \ce{H} à céder \emph{et} peut
        capter un \ce{H+}. \ce{H3PO4} (acide pur) et \ce{PO4^{3-}} (base pure) ne sont pas des ampholytes.
  \item \ce{HPO4^{2-}} comme \textcolor{acidecol}{\textbf{acide}} : $\ce{HPO4^{2-} + H2O <=> PO4^{3-} + H3O+}$.\\
        \ce{HPO4^{2-}} comme \textcolor{basecol}{\textbf{base}} : $\ce{HPO4^{2-} + H2O <=> H2PO4- + OH-}$
        (ou, de façon équivalente, $\ce{HPO4^{2-} + H3O+ <=> H2PO4- + H2O}$).
\end{enumerate}
\end{corrbox}

\begin{corrbox}[la glycine, un acide aminé amphotère]
\begin{enumerate}[label=\alph*),leftmargin=1.6em,itemsep=3pt]
  \item Le carboxyle \ce{-COOH} perd son \ce{H+} et devient \ce{-COO-} :
  \[ \ce{HOOC-CH2-NH3+} \;\longrightarrow\; \ce{^{-}OOC-CH2-NH3+} + \ce{H+}. \]
  L'espèce obtenue porte une charge $-$ (carboxylate) \emph{et} une charge $+$ (ammonium) : charge globale
  nulle, c'est le \textbf{zwitterion}.
  \item Le groupe ammonium \ce{-NH3+} perd à son tour son \ce{H+} et devient l'amine \ce{-NH2} :
  \[ \ce{^{-}OOC-CH2-NH3+} \;\longrightarrow\; \ce{^{-}OOC-CH2-NH2} + \ce{H+}. \]
  \item Le zwitterion est un \textbf{ampholyte} :
  \begin{itemize}[leftmargin=1.4em,topsep=1pt,itemsep=1pt]
     \item par son groupe \textcolor{acidecol}{ammonium \ce{-NH3+}}, il peut encore \emph{céder} un \ce{H+}
           (rôle d'\textbf{acide}), couple \ce{^{-}OOC-CH2-NH3+}/\ce{^{-}OOC-CH2-NH2} ;
     \item par son groupe \textcolor{basecol}{carboxylate \ce{-COO-}}, il peut \emph{capter} un \ce{H+}
           (rôle de \textbf{base}), couple \ce{HOOC-CH2-NH3+}/\ce{^{-}OOC-CH2-NH3+}.
  \end{itemize}
  Jouant les deux rôles, le zwitterion est bien amphotère.
\end{enumerate}
\end{corrbox}

\begin{corrbox}[lire une table et raisonner sur les couples]
\begin{enumerate}[label=\alph*),leftmargin=1.6em,itemsep=3pt]
  \item L'espèce \textbf{\ce{HC2O4-}} apparaît comme base conjuguée (de \ce{H2C2O4}) \emph{et} comme acide (du
        couple \ce{HC2O4-}/\ce{C2O4^{2-}}). Cela prouve que \ce{HC2O4-} est un \textbf{ampholyte}.
  \item Base conjuguée de \ce{CH3COOH} : \ce{CH3COO-}. \quad Acide conjugué de \ce{NH3} : \ce{NH4+}.
  \item \textbf{Non}, ils n'appartiennent pas au même couple. On passe de \ce{H2C2O4} à \ce{C2O4^{2-}} en
        retirant \emph{deux} \ce{H+} (via l'intermédiaire \ce{HC2O4-}). Un couple conjugué n'est séparé que par
        \emph{un seul} \ce{H+} ; ici il y a deux couples successifs (\ce{H2C2O4}/\ce{HC2O4-} puis
        \ce{HC2O4-}/\ce{C2O4^{2-}}).
\end{enumerate}
\end{corrbox}

\begin{astuce}
\textbf{Ce que le concours teste ici.} Reconnaître un ampholyte au fait qu'il figure des \emph{deux} côtés
d'une table de couples (comme \ce{HC2O4-} ou \ce{HCO3-}), et ne jamais confondre « deux espèces séparées par
deux protons » avec « un couple conjugué » (séparé par un seul \ce{H+}).
\end{astuce}

\end{document}
